\chapter{What a Resolution Would Mean}
\label{ch:consequences}

Suppose the question were settled. This chapter describes the two worlds, the one in which $\Pc = \NP$ and the one in which it does not, and what each would mean for mathematics, for cryptography, for optimisation, and for the Monte Carlo program of this book. It also describes the intermediate possibilities, which are many, and which are where most of the field expects the truth to lie.

\section{If P = NP}

If $\Pc = \NP$ by an algorithm with a small polynomial, the consequences are enormous and have been described by Aaronson, Fortnow and others at length. Every problem in the polynomial hierarchy is in $\Pc$. Every optimisation problem whose solutions can be evaluated quickly can be solved to optimality quickly, since the decision version is in $\NP$ and the optimum is found by binary search. Cryptography based on computational hardness is impossible: a one way function would give a problem in $\NP$ not in $\Pc$. Public key encryption, digital signatures and the whole infrastructure of secure communication would fall, and only information theoretic schemes such as the one time pad would survive.

Mathematics itself changes. By G\"odel\index{G\"odel, Kurt}'s observation, the search for a proof of length $n$ of any statement is in $\NP$, so a polynomial time algorithm would find proofs of length $n$ in time polynomial in $n$. Every theorem with a proof of reasonable length would be found by machine, and the role of the mathematician would be to choose which statements to ask about. Learning becomes easy in a precise sense: finding the smallest circuit consistent with given data is in $\NP$, so Occam's razor could be implemented, and machine learning would reduce to the choice of a representation. Levin's universal search algorithm would find a witness for every yes instance of an $\NP$ problem in time within a constant of optimal, without anyone knowing what the optimal algorithm is.

If $\Pc = \NP$ by an algorithm with running time $n^{1000}$, or by a nonconstructive argument that shows an algorithm exists without exhibiting it, the practical consequences are nil and the theoretical consequences are a reorganisation of the subject around the question of the exponent. Few expect this, but Levin's algorithm shows that constructiveness is not the issue: if $\Pc = \NP$ then a specific algorithm, known today, finds a satisfying assignment of every satisfiable formula in polynomial time; only the polynomial is unknown.

For the Monte Carlo program, $\Pc = \NP$ makes everything easy and nothing interesting. The queens level sets could be sampled almost uniformly by the self reducibility\index{self reducibility} route from a polynomial time approximate counter, the ladder would be unnecessary, and the thesis that hardness lives in the widths of a staircase would be true but vacuous, since nothing would be hard.

\section{If P is not NP}

If $\Pc \neq \NP$, which is what everyone expects, then the landscape is as it appears: $\NP$ complete problems have no polynomial time algorithm, the polynomial hierarchy is infinite or at least does not collapse to $\Pc$, and whether cryptography is possible turns on whether one way functions exist. A proof would be a landmark but would change nothing in practice, since practice already assumes it. What would change is the theory of lower bounds, which would at last have a superpolynomial lower bound against uniform polynomial time for an explicit problem in $\NP$, and the technique that gave it would likely give more.

For the Monte Carlo program, $\Pc \neq \NP$ is consistent with both $\NP = \RP$ and $\NP \neq \RP$, and so says nothing about whether the queens chain mixes. What rules out the possibility that every conditioned board mixes rapidly with inverse polynomial rung ratios is $\NP \not\subseteq \BPP$, which is expected to hold: by the hardness versus randomness theorems, if $\NP \subseteq \BPP$ then under the circuit lower bound hypothesis $\NP = \Pc$, so $\NP \not\subseteq \BPP$ follows from $\Pc \neq \NP$ together with the hypothesis. In the expected world the queens chain is expected to mix torpid\index{torpid mixing}ly at low energy, and the program's Monte Carlo half fails for the reason that it should.

\section{The intermediate worlds}

Impagliazzo\index{Impagliazzo, Russell}'s five worlds of Chapter~\ref{ch:history} are the standard map of the possibilities between the two extremes, and there are finer distinctions.

$\NP$ hard on the worst case but easy on average. In Heuristica, every $\NP$ problem has a polynomial time algorithm that succeeds on all but a negligible fraction of inputs under every samplable distribution, but some inputs are hard. For the queens this would mean that a random partial placement is easy to complete or to show incomplete, but that specific adversarial placements, such as those produced by the Gent Jefferson Nightingale reduction from a hard formula, are not. The Monte Carlo walk would then succeed on random instances and fail on constructed ones, which is in fact what the experimental evidence on random constraint satisfaction suggests happens for satisfiability away from the threshold.

$\NP$ hard on average but no one way functions. In Pessiland hard instances can be generated but not with a planted solution, so cryptography is impossible and so is the efficient generation of hard instances with known answers. For counting this is the world in which the queens are hard to count on random boards but there is no way to make a board whose count one knows.

$\Pc \neq \NP$ but $\NP \subseteq \BPP$. This is the world in which randomness is a genuine resource, the Impagliazzo Wigderson hypothesis fails, and some polynomial time sampler for the queens level sets exists after all. Nobody believes in it, and the belief is well founded, but it has not been ruled out, and it is the world in which the Monte Carlo program would succeed in full.

$\NP \subseteq \Ppoly$ but $\Pc \neq \NP$. Here $\NP$ problems have small circuits that cannot be found uniformly. Karp\index{Karp, Richard} and Lipton showed this collapses the hierarchy to the second level, which is believed not to happen, and so this world is disbelieved for a reason stronger than taste.

\section{The role of counting}

Whatever the world, counting sits above deciding. Toda's theorem places the whole hierarchy below $\Pc^{\sharpP}$ unconditionally, so even if $\Pc = \NP$ exact counting is harder than everything in the hierarchy unless the hierarchy collapses all the way to $\Pc^{\sharpP}$, which would require $\sharpP$ functions to be computable in polynomial time, which is $\Pc = \mathrm{PP}$, a collapse stronger than $\Pc = \NP$. Approximate counting is more sensitive: by Stockmeyer\index{Stockmeyer's theorem} it is in $\BPP^{\NP}$, so if $\Pc = \NP$ it is in $\Pc$. The vertical walk is an approximate counting algorithm, and so its complexity tracks $\NP$ rather than $\sharpP$: it is polynomial if and only if the level sets can be sampled, which for the queens, whose completion problem is $\NP$ complete, is if and only if $\NP = \RP$.

\section{The program in each world}

We may summarise. In Algorithmica the program is trivial. In Heuristica it succeeds on random boards and fails on constructed ones, and the failure mode would be a chain trapped in a component of a level set that stays shattered to a depth growing with $N$, so that no rung the walk can reach has a usable ratio to the ground state. In Pessiland, Minicrypt and Cryptomania (defined in Chapter~\ref{ch:history}), which differ in matters of cryptography that the program does not touch, the program fails on hard instances for the same reason. In the disbelieved world where randomness beats determinism, the program succeeds and the hardness versus randomness paradigm fails. In every world the staircase is the same and the question is the same: can the bottom rungs be reached? The answer is the answer to $\NP = \RP$, which in every expected world is no.

What the program contributes is not an answer but a place to look. The question whether $\NP = \RP$ has been, since Jerrum, Valiant\index{Valiant, Leslie} and Vazirani\index{Jerrum, Valiant and Vazirani}, a question about the mixing of chains on the level sets of $\NP$ complete energies, and nobody has found a chain that mixes or proved that none does for a natural energy. The queens are as natural an energy as there is. Part IV sets up the program in full, gives the best Monte Carlo analysis we can, and leaves the chain where we found it: slow at the bottom, for reasons we can see and cannot yet prove.
